\uuid{ClMr}
\titre{Intégration par parties}
\auteur{hanani ; mzari}
\organisation{exo7}
\licence{CC BY-NC-SA 4.0}
\datecreate{2020-01-20}
\niveau{L1}
\module{Analyse}
\chapitre{Calcul d'intégrales}
\sousChapitre{Primitives diverses}
\difficulte{3}
\theme{Primitives, intégration par parties}
\exo7qcmid{lille-2.primitives.035}

\contenu{
\begin{question}
\qclassification{127.11}
\qtags{motcle=intégration par parties}
On se place sur $]-1,+\infty[$. Parmi les affirmations suivantes, cocher celles qui sont vraies :
\begin{answers}
\bad{$\displaystyle \int \ln (x+1)\, \mathrm{d}x=\mathrm{e}^{x+1}+k$, où $k\in \Rr$.}
\good{$\displaystyle \int \ln (x+1)\, \mathrm{d}x=(x+1)\ln (x+1)-\int \mathrm{d}x$.}
\bad{$\displaystyle \int x\ln (1+x)\mathrm{d}x=\int x\mathrm{d}x\times\ln (1+x)=\frac{x^2}{2}\ln (1+x)+k$, où $k\in \Rr$.}
\good{$\displaystyle \int x\ln (1+x)\mathrm{d}x=\frac{x^2}{2}\ln (1+x)-\int \frac{x^2\mathrm{d}x}{2(1+x)}$.}
\end{answers}
\begin{explanations}
Une intégration par parties avec $u=\ln (x+1)$ et $v=x+1$ donne
$$\int \ln (x+1)\, \mathrm{d}x=(x+1)\ln (x+1)-\int \mathrm{d}x=(x+1)\ln (x+1)-x+k,\; k\in \Rr.$$
Une intégration par parties avec $u=\ln (x+1)$ et $v=x^2/2$ donne
$$\int x\ln (1+x)\mathrm{d}x=\frac{x^2}{2}\ln (1+x)-\int \frac{x^2\mathrm{d}x}{2(1+x)}.$$
\end{explanations}
\end{question}
}
